The multiplicities of the two roots of $P$ are either $(2,2)$ or $(3,1)$.

In the first case write $P(x)=(x-r)^2(x-s)^2$ and $d=s-r\ne0$. Its only critical point away from its roots is $\dfrac{(r+s)}{2}$, where $P=\dfrac{d^4}{16}$. The polynomial $P-1$ has exactly three distinct roots precisely when this simple critical point has value $1$, so $d^4=16$. Since one root is $3$, the four distinct polynomials are $(x-3)^2(x-3-d)^2$ as $d$ runs through the fourth roots of $16$. Their values at $1$ are $(4+2d)^2$, whose sum is $64$.

In the second case write $P(x)=(x-r)^3(x-s)$, with $d=s-r$. Its nonroot critical point is $\dfrac{(r+3s)}{4}$, with value $-27\dfrac{d^4}{256}$, so $d^4=-\dfrac{256}{27}$. If the triple root is $3$, then $P(1)=16+8d$, summing to $64$ over the four choices of $d$. If the simple root is $3$, then $P(1)=-2(d-2)^3$, also summing to $64$, since the sums of the first three powers of those four $d$ values vanish.

These twelve polynomials exhaust the cases, and all their displayed critical points are simple, so each indeed has exactly three roots at level $1$. The total is $\boxed{192}$.
